%!TEX root = ../main.tex
\section{Related Work}
\label{sec:related}

\paragraph{Foundation models for paralinguistic audio.} Frozen self-supervised speech models (Wav2Vec2~\cite{baevski2020wav2vec2}, HuBERT~\cite{hsu2021hubert}, WavLM~\cite{chen2022wavlm}) have become the standard substrate for speaker-, emotion-, and health-related audio classification. Pasad~\etal~\cite{pasad2021layer} and Chen~\etal~\cite{chen2022wavlm} document a layer-wise functional stratification in these models: early Transformer layers concentrate speaker identity and acoustic detail, mid-to-late layers concentrate phonetic and paralinguistic content. This stratification motivates layer-weighted pooling for downstream heads. We use it twice: A2's softmax layer-weight-fused pooled stats (the canonical foundation-model head architecture for this corpus), and A2.5's per-layer cold/speaker audit (Section~\ref{sec:method}) that re-initialises the layer-weight softmax with a data-derived prior.

\paragraph{The speaker shortcut on small-data paralinguistic corpora.} Huckvale's analysis of the 2017 Cold sub-challenge submissions~\cite{huckvale2017cold} surfaced what subsequent work calls the speaker shortcut: classifiers trained on representations that retain speaker-identifying information can score on the paralinguistic task by indirectly recognising speakers, especially when the train/test split is not strictly speaker-disjoint. The 2017 ComParE Cold sub-challenge baseline was 0.710 UAR on the hidden test set with a 6373-d ComParE-2016 acoustic feature set~\cite{schuller2017interspeech}. A range of submissions used speaker-grouped cross-validation but few systematically audited what their probes were measuring, motivating the speaker probe protocol we use here (Section~\ref{sec:bg}).

\paragraph{De-confounding mechanisms.} Three families of de-confounding interventions are commonly proposed in the audio paralinguistic literature.

\textit{Data-level} interventions reshape the training distribution to break the cold-vs-speaker correlation. Cross-speaker audio splicing, gain/noise augmentation, and prosody-preserving voice conversion are typical. We test cross-speaker splicing and an embedding-mixup variant (A5.5 in Section~\ref{sec:deconf}); both fail for distinct mechanism reasons documented in M8 and M9.

\textit{Representation-level} interventions reshape the embedding geometry, typically with supervised contrastive losses~\cite{khosla2020supcon}. A common variant masks speaker-positive pairs to encourage same-class-different-speaker proximity. We test this (A6 in Section~\ref{sec:deconf}) and find it dominated by simpler cold-CE training plus dimensionality-bottleneck effects (M10, M11).

\textit{Gradient-level} interventions subtract the speaker direction from the optimization signal via an adversarial discriminator with gradient reversal (DANN~\cite{ganin2015dann}, MDD~\cite{zhang2019mdd}). We test DANN at two scopes (A7, A7c in Section~\ref{sec:deconf}); both fail due to substrate-resistance: the discriminator memorises individual training-chunk speaker IDs without generalising, so the adversary's gradient pushes against memorised fingerprints rather than generalising speaker structure (M12, M13, M14).

\paragraph{Audit-gated late fusion.} Logit-level late fusion of WavLM with handcrafted acoustic features is closer to the 2017 baseline tradition (kitchen-sink ComParE-2016) than to end-to-end deep models. Our A5b stage admits low-dimensional handcrafted feature groups (G1--G6, defined in Section~\ref{sec:method}) into a constrained late-fusion architecture by a per-group subtractive cold-vs-speaker score (cold gain $-$ speaker gain). The architecture is logit-additive: $\ell_{\text{fused}} = \ell_{\text{A2.5}} + \beta \cdot \mathrm{mean}_g(\mathrm{zscore}_g(\ell_g))$. The K-locked $\beta$-sweep methodology is what produces the K=2 LOCKED canonical (Section~\ref{sec:method}); the design choice of locking K rather than free-sweeping K together with $\beta$/$\tau$ is itself a methodology contribution discussed in Section~\ref{sec:disc} (the M4 $\tau$-sweep pathology).
